\documentclass[11pt]{article}
\usepackage[utf8]{inputenc}
\usepackage{amsmath,amssymb,amsthm}
\usepackage{booktabs}
\usepackage[margin=2.5cm]{geometry}
\usepackage{hyperref}

\newtheorem{conjecture}{Conjecture}
\newtheorem{observation}{Observation}
\newtheorem{proposition}{Proposition}
\theoremstyle{remark}
\newtheorem{remark}{Remark}

\newcommand{\D}{\mathrm{D}}

\title{Differential expansion for non-rectangular representations:\\
lines $A=q^{-m}$, $U_q\mathfrak{gl}(k|k+m)$ and the hook towers\\[2mm]\large\emph{(draft, continues docs/notes/defect\_tower.tex)}}
\author{}
\date{}

\begin{document}
\maketitle

\begin{abstract}
For rectangular $R=[r^s]$ the knot-independent content of the differential expansion (DE) is the
factor $\D_r\D_{-s}$ of $H_R-1$; for non-rectangular $R$ no such factor exists.  We show that its role is
played by the restrictions of $H_R$ to the lines $A=q^{-m}$ ($\D_m=0$), where $H_R$ is a
$U_q\mathfrak{gl}(k|k+m)$ invariant for every $k\ge0$.  (i) All coincidences $H_R=H_{R'}$ on these lines are
Berezinian twists of $\mathfrak{gl}(k|k+m)$; this explains the rule found earlier and adds the rules with $k\ge2$
(e.g.\ $H_{[2,2,2]}=H_{[3,3]}$ and $H_{[3,2,2]}=H_{[4,3]}$ at $A=1$), with no exception on the knot table ($|R|\le6$) and on double braids
($|R|\le10$).  (ii) Along every family $R=[r,\mu]$ with $\mu_1\le m$ the restriction is a Laurent polynomial in
$\lambda=q^r$ (a typical $\mathfrak{gl}(1|m+1)$ module of dimension $2^{m+1}\dim\mu$), whose $\lambda$-degree
$e^\mu_m$ generalises the tower $e_m$ of the symmetric DE.  (iii) The hook family $[r,1]$ detects the Seifert genus where the
symmetric family does not: $e^{[1]}_1=2g$ for every knot we tested (216 knots), including all 36 knots with
$\deg\Delta<g$ up to 12 crossings, 27 of which have $e_1<2g$, and the mutant pairs with different genera
(Conway / Kinoshita--Terasaka: $6$ vs $4$); for $m=2$, $e^{[1]}_2=3g$ where computed.  The exact symmetric $e_1$
also splits the undecided genus-3 knots: $e_1=5$ (slope $5/2$ if it persists) for $12n_{256},12n_{257},12n_{264},
12n_{267},12n_{665},12n_{812}$ and $e_1=4$ for the others.  Hence the defect of the DE for generic $R$ is governed by the
genus, and the symmetric defect is only its lower envelope.
\end{abstract}

\section{Lines and supergroups}

Conventions as in \texttt{defect\_tower.tex}: $H_R(A,q)$ reduced, topological framing, $\D_n=Aq^n-A^{-1}q^{-n}$,
$H_{R^T}(A,q)=H_R(A,-q^{-1})$; the line $A=q^{-m}$ is $\D_m=0$, and $R\to R^T$ maps the line $m$ to $-m$.

\paragraph{Dictionary.} The Hecke algebra acts on $(\mathbb{C}^{k|n})^{\otimes N}$ through the Perk--Schultz
matrix, and its Markov trace with $A=q^{k-n}$ is the quantum supertrace (super Schur--Weyl duality).  Covariant
irreducible $\mathfrak{gl}(k|n)$-modules $V_R$ are labelled by the diagrams in the $(k|n)$-hook,
$R_{k+1}\le n$, with highest weight
\begin{equation}\label{hw}
  \big(R_1,\dots,R_k\ \big|\ \langle R^T_1-k\rangle,\dots,\langle R^T_n-k\rangle\big),\qquad \langle x\rangle=\max(x,0).
\end{equation}
With $n=k+m$ the $(1,1)$-tangle of the $R$-cabled knot acts on $V_R$ by $H_R(A=q^{-m},q)$ (mirror image;
the argument of Proposition~1 of \texttt{defect\_tower.tex}, which never used $k=1$).  Every $R$ lies in the
$(k|k+m)$-hook for all large $k$, so \emph{every} $H_R$ restricted to \emph{every} line is a supergroup invariant.

We checked the dictionary independently for $k=1$ with a vertex model of $U_q\mathfrak{gl}(1|m+1)$ built only from
the Perk--Schultz matrix: $V_R$ is cut out of $V_{R-\square}\otimes V$ as the eigenspace of the monodromy with
eigenvalue $q^{2c(\square)}$, and all braidings are fused recursively, exactly modulo a prime
(\texttt{scripts/super\_vertex.py}).  It reproduces $H_R(q^{-1},q)$ of the table for all $R$ with
$R_2\le2$, $|R|\le6$, on the knots tested ($3_1,4_1,5_2,6_2,8_{20}$; $11n_{34},11n_{42}$ for $|R|\le4$),
including the dimensions: $[r]\mapsto4$, $[r,1]\mapsto8$, $[r,1,1]\mapsto12$, and $[r,2]\mapsto4$.

\section{Universal reductions are Berezinian twists}

The Berezinian $\mathrm{Ber}$ of $\mathfrak{gl}(k|k+m)$ has weight $(1,\dots,1|-1,\dots,-1)$, and
$V_R\otimes\mathrm{Ber}^{\pm1}$ is again covariant whenever the weight stays in the hook.  A one-dimensional twist
changes only the framing factor, hence:

\begin{proposition}[rule $b_k$]
On the line $A=q^{-m}$, for every $k\ge1$ with $m+k\ge0$:
\[
  R_{k+1}=m+k\quad\Longrightarrow\quad H_R=H_{R'},\qquad R'=(R_1+1,\dots,R_k+1,R_{k+2},R_{k+3},\dots).
\]
\end{proposition}
\begin{itemize}
\item $k=1$ is rule (b) of \texttt{defect\_tower.tex}: $R_2=m+1\Rightarrow R\equiv(R_1+1,R_3,\dots)$, the second row
slides into the first.  For $m=0$ it is the hook property $H_R(A=1)=\Delta(q^{2|R|})$ ($\mathfrak{gl}(1|1)$).
\item $m+k=0$ ($\mathfrak{gl}(k|0)$) is $\mathfrak{sl}_k$: a column of height $k$ can be added.
\item $k=2$ is new: at $A=1$, $[2,2,2]\equiv[3,3]$ (knot table) and $[3,2,2]\equiv[4,3]$, $[4,2,2]\equiv[5,3]$,
$[5,2,2]\equiv[6,3]$, $[6,2,2]\equiv[7,3]$ (double braids); predicted further: $[2,2,2,1]\equiv[3,3,1]$,
$[3,3,2]\equiv[4,4]\equiv[2,2,2,2]$ at $A=1$, $[3,3,3]\equiv[4,4]$ at $A=q^{-1}$.
\end{itemize}
Together with the same rules for $R^T$ on the line $-m$ and $\mathfrak{sl}$ duality (rule (a)), these generate an
equivalence relation on diagrams for every $m$ (\texttt{scripts/super\_reductions.py}).

\begin{observation}
The Berezinian classes are exactly the coincidence classes of colored HOMFLY on the lines:
\begin{itemize}
\item knot table, 60 random knots, all $|R|\le6$, $-8\le m\le8$: 206 predicted identities, all hold, no other
identity holds;
\item antiparallel double braids $H(1,\pm1),H(2,\pm1),H(2,\pm2)$ ($3_1,4_1,5_2,6_1,7_4,8_3$) by the strengthened
Formula II, 57 diagrams up to $|R|=10$, $m\in\{-2,\dots,3\}$: 241 predicted identities, all hold, no other.
\end{itemize}
\end{observation}
\begin{remark}
\texttt{defect\_tower.tex} states that (a), (b) are complete for $|R|\le6$.  This is not so: $[2,2,2]\equiv[3,3]$ at $A=1$
holds for every knot.  The check in \texttt{scripts/line\_reductions.py} compared normal forms obtained by
iterating the reduction in one direction only, so identities reached through a common descendant were counted as
``predicted'' and $(b_2)$ was never tested.
\end{remark}

\paragraph{Meaning for the DE.}  For rectangular $R=[r^s]$ the class of $R$ contains $\emptyset$ on the lines
$\D_r=0$ and $\D_{-s}=0$, which is the factor $\D_r\D_{-s}$ of $H_R-1$.  A non-rectangular $R$ is never
equivalent to $\emptyset$ (no factor of $H_R-1$), but on every line it is equivalent to other diagrams, so
$H_R-H_{R'}$ is divisible by $\D_m$ whenever $R\sim R'$ on $\D_m=0$.  This is the knot-independent part of the
non-rectangular DE.  In the channel form $H_R=1+\sum_cZ_R^cG_c$ it requires
$\sum_c(Z_R^c-Z_{R'}^c)G_c\equiv0\pmod{\D_m}$ for every knot.  For the factors $Z_R^c=F_c(R)$ of Formula~II this does
\emph{not} hold channel by channel: for $|R|\le6$, $m=1,2$, $F_c(R)\equiv F_c(R')$ on the line for 25 of the 35
equivalent pairs, and in the other 10 the difference is confined to blocks formed by a single $\lambda$ and the pairs
$\{Z,Z'\}$ with $Z\cup Z'\subseteq\lambda$ (e.g.\ $[2,1]$ with $\{[2],[1,1]\}$, $[3,1]$ with $\{[3],[2,1]\}$).
So on the lines the knot functions of a single and of the pairs below it are tied by linear relations, the
non-rectangular counterpart of ``$\D_m$ divides $G_j$''.

\section{Families on a line: the towers $e^\mu_m$}

Fix $m\ge1$.  By (\ref{hw}) with $k=1$, $R=[r,\mu]$ with $\mu_1\le m+1$ is the $\mathfrak{gl}(1|m+1)$-module with
weight $(r\,|\,\nu)$, $\nu=(\mu^T_1,\dots,\mu^T_{m+1})$; modulo $\mathrm{Ber}$ one may take $\mu_1\le m$.  For large $r$
it is a typical Kac module of dimension $2^{m+1}\dim_{\mathfrak{gl}_{m+1}}\nu$, and $r$ enters only through the
continuous weight $q^r$.  Hence

\begin{observation}
For every knot, $m\ge1$ and $\mu$ with $\mu_1\le m$,
\[
  P^\mu_m(q,\lambda):=H_{[r,\mu]}(A=q^{-m},q)\big|_{q^r=\lambda}
\]
is a Laurent polynomial in $\lambda$ (and $q$), symmetric under $\lambda\to\lambda^{-1}$ in the normalisation of
\texttt{super\_vertex.py}, with $P^\mu_m(1,\lambda)=\Delta(\lambda^2)^{m+1}$.
\end{observation}
Checks: in the vertex model ($r\le30$--$40$) for $\mu=\emptyset,[1]$ at $m=1$ on 214 knots, $\mu=[1,1],[1,1,1]$ on
$3_1,4_1,6_2$, and $\mu=\emptyset,[1]$ at $m=2$ on the knots of the table below; the exponents of $\lambda$ are read off by
Berlekamp--Massey from the sequence in $r$, so polynomiality is tested, not assumed (otherwise the recurrence would not
stabilise), and the exponent sets are always symmetric, $\{-2e,\dots,2e\}$ in steps of 2 (the coefficients were
compared for $3_1$ and $11n_{34}$);
the double braids up to $|R|=10$ fit $\lambda$-degree $2(m+1)$ with independent extra points for $m=1$
($\mu=\emptyset,[1],[1,1],[1,1,1]$) and $m=2$ ($\mu=\emptyset,[1]$), 42 fits, none fails;
$P^{[1]}_1(1,\lambda)=\Delta(\lambda^2)^2$ exactly for $3_1$ and for $11n_{34}$ (where $\Delta=1$ but the
$\lambda$-span is 24).  $\mu=\emptyset$ is the tower $P_m$ of \texttt{defect\_tower.tex}.

\noindent Define $e^\mu_m:=\tfrac12\max\deg_\lambda P^\mu_m$, so that $e^\emptyset_m=e_m$.  For $m=0$
all hooks give $\Delta(q^{2|R|})$, so $e^\mu_0=\deg\Delta$ for every hook family.

\section{The hook family detects the genus}

\begin{observation}\label{obs:hook}
$e^{[1]}_1=2g$ for every knot computed: all 36 knots with $\deg\Delta<g$ up to 12 crossings, the 7 mutant partners of genus 2 below, and a control set of
173 knots with at most 10 crossings and braid index $\le4$.  Where $m=2$ was computed, $e^{[1]}_2=3g$.
\begin{center}
\begin{tabular}{lcc|cc|cc}
\toprule
knot & $\deg\Delta$ & $g$ & $e_1$ & $e^{[1]}_1$ & $e_2$ & $e^{[1]}_2$\\
\midrule
$11n_{42}$ & 0 & 2 & 3 & 4 & 5 & 6\\
$11n_{67}$ & 1 & 2 & 3 & 4 &  & \\
$11n_{97}$ & 1 & 2 & 3 & 4 &  & \\
$12n_{268}$ & 1 & 2 & 3 & 4 &  & \\
$12n_{313}$ & 0 & 2 & 3 & 4 &  & \\
$12n_{430}$ & 0 & 2 & 3 & 4 &  & \\
$12n_{51}$ & 1 & 2 & 3 & 4 &  & \\
$12n_{124}$ & 1 & 2 & 4 & 4 &  & \\
$12n_{23}$ & 1 & 2 & 4 & 4 &  & \\
$12n_{293}$ & 1 & 2 & 4 & 4 & 6 & 6\\
$12n_{321}$ & 1 & 2 & 4 & 4 &  & \\
$12n_{411}$ & 1 & 2 & 4 & 4 &  & \\
$12n_{457}$ & 1 & 2 & 4 & 4 &  & \\
$12n_{519}$ & 1 & 2 & 4 & 4 &  & \\
$11n_{34}$ & 0 & 3 & 3 & 6 & 5 & 9\\
$11n_{152}$ & 2 & 3 & 4 & 6 &  & \\
$11n_{45}$ & 2 & 3 & 4 & 6 &  & \\
$11n_{73}$ & 2 & 3 & 4 & 6 & 6 & 9\\
$12n_{129}$ & 1 & 3 & 4 & 6 &  & \\
$12n_{132}$ & 2 & 3 & 4 & 6 &  & \\
$12n_{221}$ & 2 & 3 & 4 & 6 &  & \\
$12n_{231}$ & 2 & 3 & 4 & 6 &  & \\
$12n_{28}$ & 2 & 3 & 4 & 6 &  & \\
$12n_{31}$ & 1 & 3 & 4 & 6 &  & \\
$12n_{56}$ & 2 & 3 & 4 & 6 &  & \\
$12n_{63}$ & 2 & 3 & 4 & 6 &  & \\
$12n_{808}$ & 2 & 3 & 4 & 6 &  & \\
$12n_{87}$ & 2 & 3 & 4 & 6 &  & \\
$12n_{256}$ & 2 & 3 & 5 & 6 &  & \\
$12n_{257}$ & 2 & 3 & 5 & 6 &  & \\
$12n_{264}$ & 2 & 3 & 5 & 6 &  & \\
$12n_{267}$ & 2 & 3 & 5 & 6 &  & \\
$12n_{665}$ & 2 & 3 & 5 & 6 &  & \\
$12n_{812}$ & 2 & 3 & 5 & 6 &  & \\
$12n_{750}$ & 2 & 3 & 6 & 6 & 9 & 9\\
$12n_{830}$ & 2 & 3 & 6 & 6 & 9 & 9\\
$11n_{151}$$^*$ & 2 & 2 & 4 & 4 &  & \\
$11n_{39}$$^*$ & 2 & 2 & 4 & 4 &  & \\
$11n_{74}$$^*$ & 2 & 2 & 4 & 4 &  & \\
$12n_{225}$$^*$ & 2 & 2 & 4 & 4 &  & \\
$12n_{232}$$^*$ & 2 & 2 & 4 & 4 &  & \\
$12n_{34}$$^*$ & 2 & 2 & 4 & 4 &  & \\
$12n_{57}$$^*$ & 2 & 2 & 4 & 4 &  & \\
\bottomrule
\end{tabular}
\end{center}
$e_m=e^\emptyset_m$ is the symmetric tower, now computed exactly by the vertex model for all these knots
($r\le30$; the table of \texttt{defect\_tower.tex} had only lower bounds for most of them).  $^*$: genus-2 mutant partners
of $11n_{45},11n_{73},11n_{152},12n_{28},12n_{63},12n_{231},12n_{56}$ with identical symmetric $H_{[r]}$.
\end{observation}


For $11n_{34}$ the top coefficient is
\[
  [\lambda^{\pm12}]\,P^{[1]}_1=\frac{(q^2-1)^3(q^4+q^2+1)(q^8+q^2-1)(q^8-q^6-1)}{q^{16}}
  =\{q\}^3\,[3]\,\frac{(q^8+q^2-1)(q^8-q^6-1)}{q^{11}}
\]
(vertex-model variable $q$, i.e.\ the mirror; $P^{[1]}_1(1,\lambda)=1=\Delta^2$, every $\lambda^{j}$, $j\neq0$, vanishes at $q=1$),
it vanishes at $q=1$ to third order, so all the genus information is in the $\hbar$-corrections, as for the
anomalous symmetric towers.

\paragraph{Consequences.}
\begin{itemize}
\item The answer to Question~1.4 of L\'opez Neumann--Garoufalidis (does the maximum over typical modules attain the
genus bound?) is positive on all examples, and already the 8-dimensional Kac module $(r|1,0)$ of
$\mathfrak{gl}(1|2)$ suffices where the 4-dimensional $\mathrm{Sym}^r$ fails.
\item $P^{[1]}_1$ is not a mutation invariant: it distinguishes $11n_{34}$ from $11n_{42}$ by its $\lambda$-degree
(genus 3 vs 2), and likewise the other mutant pairs with different genera below.  This agrees with the mutant
analysis of the paper: the difference sits in the pair channels of $R\otimes\bar R$, which first appear for $[2,1]$.
\item In the DE language: on $\D_1=0$ the symmetric knot functions $G_{[j]}$ vanish for $j>e_1$, but the functions of
the non-symmetric channels contributing to $[r,1]$ (singles $[j,1]$ and pairs $\{[j],[j-1,1]\}$) survive up to level
$2g$.  The natural defect of the DE for generic $R$ is therefore
\[
  1+\delta_{\rm gen}:=\lim_{m\to\infty}\frac{\max_\mu e^\mu_m}{m+1}\overset{?}{=}g,
\]
while the symmetric defect $1+\delta=\lim e_m/(m+1)\in[\deg\Delta,g]$ is its lower envelope.
\end{itemize}

\section{Open questions}
\begin{enumerate}
\item Is $e^{[1]}_m=(m+1)g$ for all $m$ (checked for $m=1$, and for $m=2$ on $11n_{34},11n_{42},11n_{73},12n_{293},12n_{750},12n_{830}$), or does one need larger $\mu$ for larger
$m$?
\item Families with $k\ge2$ (two growing rows): the even part $\mathfrak{gl}(k)$ is then finite-dimensional with
$r$-dependent dimension, so polynomiality in $q^{R_i}$ is not expected; what replaces it?
\item A closed form for the knot-independent factors adapted to the lines, i.e.\ a triangular change of the Formula~II
basis in each block (single $\lambda$, pairs below it) that makes $Z_R^c$ respect the Berezinian classes.
\end{enumerate}

\end{document}
